\begin{textsample}{POS Dim 2 – 2021\_c4 – Score 20.00 – 2021\_c4/enade\_2021\_tec\_an\_sist\_0e41d668fe.txt}  \label{ex:f2_pos_032}
CAPÍTULO 4 PERCEPÇÃO DA PROVA Nas análises feitas neste capítulo, constam as percepções dos concluintes da Área de Tecnologia em Análise e Desenvolvimento de Sistemas sobre a prova aplicada no Enade/2021.
Essas percepções foram mensuradas por meio de nove questões que avaliaram desde o grau de dificuldade da prova até o tempo gasto para concluí -la.
As percepções sobre a prova foram relacionadas ao desempenho dos estudantes, à Grande Região de funcionamento do curso, à Categoria Administrativa e ao tipo de Organização Acadêmica da IES.
O \textbf{questionário} de percepção da prova encontra -se ao final do Anexo VIII, onde está a \textbf{reprodução} do exame.
O desempenho dos estudantes foi classificado em quatro quartos.
Para tanto, esse desempenho foi \textbf{ordenado} de forma ascendente.
O percentil 25, P25, também conhecido como primeiro quartil, é a nota de desempenho que deixa um quarto ( 25 \% ) dos valores observados abaixo e três quartos acima.
Na Figura 1, apresenta -se uma ilustração desse conceito.
O quarto inferior de desempenho é composto pelas notas abaixo do primeiro quartil.
Já o percentil 75, P75, também conhecido como terceiro quartil, é o valor para o qual há três quartos ( 75 \% ) dos dados abaixo e um quarto acima dele.
O quarto superior de desempenho é composto pelas notas iguais ou acima do terceiro quartil.
O percentil 50, P50, também conhecido como mediana, é o valor que divide as notas em dois conjuntos de igual tamanho.
O segundo quarto inclui os valores entre o primeiro quartil ( P25 ) e a mediana.
O terceiro quarto contém os valores entre a mediana ( P50 ) e o terceiro quartil ( P75 ).
Vale ressaltar que percentis, quartis e medianas são pontos que, não obrigatoriamente, pertencem ao conjunto original de dados, ao passo que os quartos são \textbf{subconjuntos} dos dados originais.
A seguir, serão apresentados gráficos com resultados selecionados, relativos às nove questões avaliadas por grupos de estudantes.
Nas barras dos gráficos, apresenta -se o percentual de alunos que assinalaram uma das opções ou a soma das porcentagens daqueles que assinalaram duas ( ou três ) delas.
Por exemplo, para as questões 1 e 2, nos gráficos, é apresentada a porcentagem total de participantes que assinalaram as opções ( D ) Difícil e ( E ) Muito difícil.
Em cada barra, foram assinalados, também, os extremos do intervalo de confiança de 95 \% como linhas verticais unidas por uma linha horizontal na forma da letra H maiúscula.
O \textbf{estimador} de um \textbf{parâmetro} com um certo nível de confiança ( por exemplo, 95,0 \% ) deve conter o \textbf{parâmetro} no intervalo de confiança em 95 \% das vezes.
Na comparação entre os \textbf{estimadores} dos \textbf{parâmetros} de duas classes de uma dada categoria ( por exemplo : Norte e Nordeste nas Grandes Regiões, de primeiro e último quarto dentro de desempenho, de IES Públicas e Privada, ou de Universidades e Faculdades ), associados aos seus respectivos intervalos de confiança, diz -se que não há diferença estatisticamente significativa entre os \textbf{parâmetros} das duas categorias quando há \textbf{interseção} entre os intervalos de confiança, e que há diferença, se os intervalos de confiança são \textbf{disjuntos} ( para maiores informações \textbf{vide} Glossário ).
Nas Tabelas do Anexo II, são apresentados os valores absolutos e a distribuição percentual18 das alternativas válidas das nove questões, segundo o mesmo \textbf{recorte} de desempenho dos alunos e da Grande Região de funcionamento do curso, de Categoria Administrativa e do tipo de Organização Acadêmica da IES. 4.1 GRAU DE DIFICULDADE DA PROVA 4.1.1 COMPONENTE DE FORMAÇÃO GERAL Ao avaliarem “ Qual o grau de dificuldade desta prova na parte de Formação Geral? ” ( Questão 1 ), 32,7 \% do grupo de inscritos e presentes optaram pelas alternativas Difícil ou Muito difícil.
Entretanto, para mais da metade dos estudantes ( 55,1 \% ), o Componente de Formação Geral da prova foi avaliado como de grau de dificuldade Médio ( Gráfico 4.1, Gráfico 4.2 e, no Anexo II, a Tabela II.1 ).
O percentual de estudantes que consideraram a prova Difícil ou Muito difícil foi maior na região Norte, onde a proporção foi de 40,7 \%, enquanto a de menor proporção foi no Nordeste, com 23,8 \%.
No Gráfico 4.1, é possível observar que as diferenças entre a região Nordeste e as demais regiões são estatisticamente significativas, assim como a diferença entre as regiões Norte e Sudeste.
Nas Grandes Regiões, a proporção de presentes à prova que avaliaram o Componente de Formação Geral como de grau de dificuldade Médio esteve entre 50,9 \%, na região Norte, e 62,9 \%, na região Nordeste. 18 Cumpre lembrar uma das convenções para tabelas \textbf{numéricas} ( ao final do \textbf{Sumário} ) sobre a possibilidade de a soma das partes não resultar em 100 \% por questões de \textbf{arredondamento}. entre as regiões Norte e Sudeste.
Nas Grandes Regiões, a proporção de presentes à prova que avaliaram o Componente de Formação Geral como de grau de dificuldade Médio esteve entre 50,9 \%, na região Norte, e 62,9 \%, na região Nordeste.
O percentual de alunos que consideraram o componente de Formação Geral da prova Difícil ou Muito difícil foi decrescente em função dos quartos de desempenho : 41,0 \%, no primeiro quarto, e 23,3 \%, no quarto, o grupo de melhor desempenho na prova.
Nos quartos de desempenho intermediários, a proporção de alunos que consideraram a prova Difícil ou Muito difícil foi de 36,2 \%, no segundo quarto, e de 30,3 \%, no terceiro quarto.
As diferenças entre todos os quartos de desempenho são estatisticamente significativas.
Para todos os quartos de desempenho, a alternativa modal a essa pergunta foi a Médio, com 48,9 \% e 61,1 \% dos respondentes nos quartos extremos, primeiro e quarto, respectivamente.
O percentual de estudantes que consideraram o componente de Formação Geral da prova Difícil ou Muito difícil foi maior entre os estudantes de Universidades, alcançando 35,0 \%, enquanto entre os estudantes das Faculdades, a Organização Acadêmica com a menor incidência, alcançou 29,9 \%.
No Gráfico 4.3, é possível observar que as diferenças entre Universidades e as demais Organizações Acadêmicas são estatisticamente significativas.
Nas Organizações Acadêmicas, a proporção de presentes à prova que avaliaram o Componente de Formação Geral como de grau de dificuldade Médio esteve entre 53,3 \%, nas Universidades, e 58,1 \% nos CEFET/IF ( ver também a Tabela II.2, no Anexo II ).
O percentual de estudantes que consideraram o componente de Formação Geral da prova Difícil ou Muito difícil foi maior entre os estudantes de cursos de IES Privadas ( 34,0 \% ), com diferença estatisticamente significativa da proporção registrada em IES Públicas ( 29,0 \% ).
Já a proporção de presentes à prova que avaliaram o Componente de Formação Geral como de grau de dificuldade Médio foi de 57,5 \%, entre os estudantes de IES Públicas, e de 54,2 \%, entre os de IES Privadas ( ver Gráfico 4.4 ). 4.1.2 COMPONENTE DE CONHECIMENTO ESPECÍFICO Ao responderem à Questão 2 – “ Qual o grau de dificuldade desta prova na parte de Componente Específico? ” – 58,2 \% do grupo de estudantes classificaram -na como Difícil ou Muito difícil.
O Componente modal de Conhecimento Específico da prova foi Difícil, para 47,1 \% dos estudantes.
Além disso, a prova foi considerada como de grau de dificuldade Médio por 37,8 \% dos alunos ( Gráfico 4.5, Gráfico 4.6, e, no Anexo II, a Tabela II.3 ).
A análise das respostas dos estudantes quanto ao grau de dificuldade do Componente de Conhecimento Específico da prova, agregado por Grande Região, mostra diferenças estatisticamente significativas entre proporções de alunos que a avaliaram como Difícil ou Muito difícil.
Tais diferenças são observadas entre a região com a maior proporção, a Norte ( 71,9 \% ), e todas as demais regiões : Nordeste ( 47,4 \% ), a menor proporção, Sudeste ( 58,1 \% ), Sul ( 60,4 \% ) e Centro-Oeste ( 58,7 \% ).
As diferenças entre a região Nordeste e todas as demais também são estatisticamente significativas.
O percentual de alunos que classificaram como grau de dificuldade como Médio, no Componente de Conhecimento Específico, variou de 26,7 \%, a 48,7 \%, nas regiões Norte e Nordeste, respectivamente.
Considerando -se a avaliação da dificuldade das questões do Componente de Conhecimento Específico da prova, de acordo com o desempenho dos estudantes, observa -se uma tendência crescente dos que consideraram a parte específica Difícil ou Muito difícil do primeiro ao terceiro quartos de desempenho, mas com a menor proporção no quarto superior ( 54,6 \% ), e a maior, no terceiro quarto ( 61,5 \% ).
As demais proporções foram de 57,0 \% e 59,5 \%, respectivamente, no primeiro e segundo quartos.
Observa -se que há diferença estatisticamente significativa dos resultados do quarto superior de desempenho e os dos segundo e terceiro quartos, assim como entre o quarto inferior e o terceiro quarto.
Já os que responderam que o grau de dificuldade das questões do Componente de Conhecimento Específico da prova foi Médio atingiram \textbf{percentagens} que variaram de 35,6 \%, no terceiro quarto, a 42,2 \%, no último.
Considerando -se a avaliação da dificuldade das questões do Componente de Conhecimento Específico da prova, o percentual de \textbf{estudantes} que as consideraram Difícil ou Muito difícil foi maior entre os estudantes de Universidades, atingindo 59,9 \% ; já entre os de Faculdades, a Organização Acadêmica com a menor incidência, essa proporção foi de 55,7 \%.
No Gráfico 4.7, é possível observar que a diferença de Universidades em relação a Faculdades e a CEFET/IF ( 56,1 \% ) é estatisticamente significativa.
Nas Organizações Acadêmicas, a proporção de presentes à prova que consideraram o Componente de Conhecimento Específico de grau de dificuldade Médio esteve entre 36,0 \%, nas Universidades, e 40,4 \%, nos CEFET/IF ( ver também Tabela II.4, no Anexo II ).
O percentual de estudantes que consideraram o Componente Específico da prova Difícil ou Muito difícil foi maior entre os estudantes de cursos de IES Privadas ( 59,1 \% ), com diferença estatisticamente significativa da proporção registrada em IES Públicas ( 55,5 \% ).
Já a proporção de presentes à prova que avaliaram esse componente como de grau de dificuldade Médio foi de 40,2 \%, entre os estudantes de IES Públicas, e de 37,0 \%, entre os de IES Privadas ( ver Gráfico 4.8 e Tabela II.4, no Anexo II ). 4.2 EXTENSÃO DA PROVA EM RELAÇÃO AO TEMPO TOTAL \textbf{Indagados} quanto à extensão da prova, em relação ao tempo total oferecido para a sua resolução ( Questão 3 ), os estudantes apontaram, com maior incidência, a alternativa que considerava a extensão Adequada, em todas as agregações consideradas ( Gráfico 4.9, Gráfico 4.10, e, no Anexo II, a Tabela II.5 ).
O percentual de alunos que responderam ser a extensão da prova Adequada foi de 58,2 \%.
Já 33,5 \% dos inscritos presentes consideraram que a prova foi Longa ou Muito longa, e 8,4 \% a avaliaram como Curta ou Muito curta.
Entre as Grandes Regiões, os que avaliaram a prova como Longa ou Muito longa, em relação ao tempo total destinado à sua resolução, atingiram uma porcentagem que variou de 32,4 \%, nas regiões Norte e Centro-Oeste, até 35,7 \%, na região Nordeste.
Não há diferença estatisticamente significativa entre as Grandes Regiões.
Considerando -se o desempenho dos alunos, nota -se ainda que 54,2 \% consideraram a extensão da prova Adequada, no quarto de desempenho superior, e 60,3 \%, no terceiro quarto.
Proporções intermediárias ocorreram no primeiro e segundo quartos : 58,7 \% e 59,5 \%, respectivamente.
No Gráfico 4.10, observa -se que as proporções de estudantes que consideraram a prova Longa ou Muito longa são bastante próximas nos três primeiros quartos de desempenho, sendo a maior proporção dessa resposta no último quarto de desempenho, 36,2 \%, e a menor, no terceiro quarto ( 32,0 \% ).
Pode -se constatar que há diferença estatisticamente significativa entre a proporção do quarto superior de desempenho e a dos demais.
Considerando -se a extensão da prova, o percentual de estudantes que avaliaram a prova como Longa ou Muito longa, em relação ao tempo total destinado à sua resolução, foi maior entre os estudantes de CEFET/IF, alcançando 36,7 \%, enquanto entre os estudantes de Faculdades, a Organização Acadêmica com a menor incidência, esse índice atingiu 32,0 \%, entre os de Universidades, o índice foi 32,9 \%, e entre os de Centros Universitários, 34,4 \%.
No Gráfico 4.11, é possível observar que há diferença estatisticamente significativa dos CEFET/IF em relação às Faculdades e às Universidades.
A proporção de presentes à prova que consideraram a sua extensão Adequada, em relação ao tempo total destinado à sua resolução, esteve entre 52,9 \%, nos CEFET/IF, e 59,7 \%, nas Faculdades ( ver também Tabela II.6, no Anexo II ).
O percentual de estudantes que consideraram a prova Longa ou Muito longa, em relação ao tempo total destinado à sua resolução, foi maior entre os estudantes de cursos de IES Públicas ( 35,2 \% ), havendo diferença estatisticamente significativa da proporção em relação aos de IES Privadas ( 32,8 \% ).
Já a proporção de presentes à prova que consideraram a extensão Adequada foi de 55,6 \%, entre os estudantes de IES Públicas, e de 59,1 \%, entre os de IES Privadas ( ver Gráfico 4.12 e Tabela II.6, no Anexo II ). 4.3 COMPREENSÃO DOS ENUNCIADOS DAS QUESTÕES 4.3.1 COMPONENTE DE FORMAÇÃO GERAL Em relação aos enunciados das questões do Componente de Formação Geral ( Questão 4 ), as opiniões foram positivas, já que 74,0 \% dos alunos avaliados consideraram todos ou a maioria dos enunciados das questões “ claros e objetivos ” ( Gráfico 4.13, Gráfico 4.14 e, no Anexo II, a Tabela II.7 ).
Na análise regional, a \textbf{percentagem} de estudantes que avaliaram que todos ou a maioria dos enunciados das questões do Componente de Formação Geral estavam “ claros e objetivos ” variou de 70,4 \%, na região Norte, a 77,4 \%, na região Centro-Oeste.
A análise das percepções dos estudantes sobre a clareza e a objetividade dos enunciados permite afirmar que todos, ou a maioria dos enunciados de questões relativas ao Componente de Formação Geral, foram considerados “ claros e objetivos ” para a maior parte dos respondentes ( maior ou igual a 70,4 \%, em todas as regiões, e maior ou igual a 65,8 \%, em todos os quartos de desempenho ).
Segundo o desempenho, observa -se que a proporção dos que emitiram essa opinião \textbf{cresce}, conforme o desempenho aumenta, com diferenças estatisticamente significativas entre todos os quartos de desempenho.
No quarto superior, a clareza e a objetividade de todos ou da maioria dos enunciados das questões foram percebidas por 80,9 \%.
A proporção dos presentes que consideraram que todos ou a maioria dos enunciados das questões do Componente de Formação Geral estavam “ claros e objetivos ” foi maior entre os estudantes de Faculdades, chegando a 75,5 \%, enquanto entre os de Universidades, a Organização Acadêmica com a menor incidência, esse índice atingiu 73,4 \% ; já entre os de Centros Universitários, o índice atingiu 73,6 \%, e entre os de CEFET/IF, 74,5 \%.
No Gráfico 4.15, é possível observar que as diferenças entre os quatro tipos de Organização Acadêmica não são estatisticamente significativas ( ver também Tabela II.8, no Anexo II ).
O percentual de estudantes que consideraram que todos ou a maioria dos enunciados das questões do Componente de Formação Geral estavam “ claros e objetivos ” foi maior entre os estudantes de cursos de IES Públicas ( 75,7 \% ), e com diferença estatisticamente significativa da proporção registrada em IES Privadas ( 73,4 \% ).
Um maior detalhamento pode ser encontrado no Gráfico 4.16 e na Tabela II.8, no Anexo II. 4.3.2 COMPONENTE DE CONHECIMENTO ESPECÍFICO Em relação aos enunciados das questões do Componente de Conhecimento Específico da prova, para 71,5 \% dos estudantes avaliados da Área de Tecnologia em Análise e Desenvolvimento de Sistemas, a clareza e a objetividade ( Questão 5 ) estavam presentes em todas ou na maioria das questões ( Gráfico 4.17 a Gráfico 4.20, e no Anexo II, as Tabelas II.9 e II.10 ).
A maioria dos estudantes de todas as Grandes Regiões brasileiras considerou “ claros e objetivos ” todos ou a maioria dos enunciados das questões do Componente de Conhecimento Específico da prova, percentual sempre maior ou igual a 66,8 \%.
As diferenças entre a região Sudeste e as regiões Sul e Centro-Oeste são estatisticamente significativas.
A proporção de estudantes que consideraram os enunciados das questões “ claros e objetivos ” apresenta uma tendência crescente em relação ao aumento de desempenho : mais elevada no quarto superior ( 79,5 \% ), se comparada ao quarto inferior de desempenho ( 64,2 \% ).
As diferenças entre todos os quartos de desempenho são estatisticamente significativas.
A proporção dos presentes que consideraram que todos ou a maioria dos enunciados das questões do Componente de Conhecimento Específico estavam “ claros e objetivos ” foi maior entre os estudantes de CEFET/IF, atingindo 73,9 \%.
Entre os estudantes de Universidades, a Organização Acadêmica com a menor incidência, essa proporção atingiu 70,1 \%, sendo intermediárias as proporções registradas em Centro Universitários ( 71,2 \% ) e em Faculdades ( 73,3 \% ).
No Gráfico 4.19, é possível observar que as diferenças de Universidades em relação a Faculdades e a CEFET/IF são estatisticamente significativas ( ver também Tabela II.10, no Anexo II ).
O percentual de estudantes que consideraram que todos ou a maioria dos enunciados das questões do Componente de Conhecimento Específico estavam “ claros e objetivos ” foi maior entre os estudantes de cursos de IES Públicas ( 74,7 \% ), com diferença estatisticamente significativa da proporção em IES Privadas ( 70,3 \% ).
Ver Gráfico 4.20 e Tabela II.10, no Anexo II, para maior detalhamento. 4.4 SUFICIÊNCIA DAS INFORMAÇÕES/INSTRUÇÕES FORNECIDAS Ao se avaliarem as informações/instruções fornecidas para a resolução das questões ( Questão 6 ), constatou -se que 81,2 \% dos respondentes da Área de Tecnologia em Análise e Desenvolvimento de Sistemas de todo o Brasil afirmaram que eram até excessivas ou suficientes em todas ou na maioria das questões ( Gráfico 4.21, Gráfico 4.22 e, no Anexo II, a Tabela II.11 ).
Quanto à distribuição de respondentes pelas Grandes Regiões, observa -se que a proporção de estudantes que consideraram as informações/instruções fornecidas até excessivas ou suficientes em todas ou na maioria das questões foi sempre superior ou igual a 74,2 \%, chegando a 83,4 \%, na região Nordeste.
Existem diferenças estatisticamente significativas entre a região Norte e as regiões Nordeste, Sul e Centro-Oeste, bem como entre as regiões Sudeste e Sul.
Levando -se em conta o desempenho dos participantes, observa -se uma tendência crescente das proporções de participantes que avaliaram as informações/instruções como até excessivas ou suficientes em todas ou na maioria das questões, havendo diferenças estatisticamente significativas entre todos os quartos de desempenho, conforme mostrado no Gráfico 4.22.
O percentual foi mais elevado no quarto superior ( 88,7 \% ), percentual superior à média nacional ( 81,2 \% ).
Já no quarto inferior, a suficiência das informações/instruções declarada como até excessiva, em todas ou na maioria das questões foi percebida por 74,5 \% dos respondentes.
Considerando -se as informações/instruções fornecidas para a resolução das questões ( Questão 6 ), a proporção que considerou que eram até excessivas ou suficientes em todas ou na maioria das questões foi maior entre os estudantes de Faculdades, \textbf{totalizando} 83,3 \%.
As Universidades foram a Organização Acadêmica em que houve a menor proporção, 79,8 \%.
No Gráfico 4.23, é possível observar que a diferença entre Faculdades e Universidades é estatisticamente significativa ( ver também Tabela II.12, no Anexo II ).
O percentual de estudantes que consideraram que as informações/instruções fornecidas para a resolução das questões eram até excessivas ou suficientes em todas ou na maioria das questões foi maior entre os estudantes de cursos de IES Públicas ( 84,6 \% ), havendo diferença estatisticamente significativa da proporção registrada em IES Privadas ( 80,0 \% ) ( ver Gráfico 4.24 e Tabela II.12, no Anexo II ). 4.5 DIFICULDADE ENCONTRADA AO RESPONDER À PROVA \textbf{Indagados} sobre as dificuldades com as quais se depararam ao responder à prova ( Questão 7 ), 20,5 \% dos estudantes apontaram o Desconhecimento do conteúdo.
Para 41,0 \%, a Forma diferente de abordagem do conteúdo foi indicada como dificuldade.
Já a Falta de motivação para fazer a prova foi a dificuldade apontada por 21,9 \% dos respondentes.
Considerando -se todo o Brasil, 8,9 \% dos respondentes afirmaram que não tiveram qualquer tipo de dificuldade para responder à prova ( Tabelas II.13 e II.14, no Anexo II ).
Nos Gráficos 4.25 a 4.28, são apresentados os percentuais de estudantes que apontaram o Desconhecimento do conteúdo como “ dificuldade ao responder à prova ”.
Na análise por Grandes Regiões, o percentual de inscritos e presentes que apontaram o Desconhecimento do conteúdo como “ dificuldade ao responder à prova ” não superou 22,7 \%.
Os percentuais variaram de 17,7 \%, na região Nordeste, a 22,7 \%, na região Norte.
A Forma diferente de abordagem do conteúdo foi a escolha modal dos estudantes, com percentuais que variaram de 39,9 \% ( região Nordeste ) a 44,9 \% ( região Norte ).
O percentual de alunos que citaram a Falta de motivação como dificuldade variou de 17,9 \% ( região Centro-Oeste ) a 24,1 \% ( região Sul ).
Os que declararam não ter qualquer tipo de dificuldade para responder à prova variaram de 6,5 \%, na região Norte, a 10,8 \%, na região Centro-Oeste.
Em relação aos quartos de desempenho, o Desconhecimento do conteúdo foi a opção escolhida por 19,9 \% dos estudantes do segundo quarto e por 20,9 \% do terceiro quarto, sem caracterizar uma tendência.
Não são observadas diferenças estatisticamente significativas entre os quartos de desempenho.
A alternativa modal entre os alunos, quando agregados pelos quartos de desempenho, para a dificuldade encontrada, foi causada pela Forma diferente de abordagem do conteúdo : 38,8 \%, no último quarto, e 42,8 \%, no segundo quarto, escolheram essa alternativa.
Na análise por Organização Acadêmica, o percentual de inscritos e presentes que apontaram o Desconhecimento do conteúdo como “ dificuldade ao responder à prova ” não superou 21,1 \%.
Os percentuais variaram de 17,8 \%, nos CEFET/IF, a 21,1 \%, nas Faculdades.
A Forma diferente de abordagem do conteúdo foi a escolha modal dos estudantes, com percentuais que variaram de 38,3 \% ( Faculdades ) a 42,5 \% ( CEFET/IF ).
O percentual de alunos que citaram a Falta de motivação como dificuldade variou de 20,9 \% ( Centros Universitários ) a 23,6 \% ( CEFET/IF ).
Os que declararam não ter qualquer tipo de dificuldade para responder à prova variaram de 6,9 \%, nos CEFET/IF, a 9,7 \%, nos Centros Universitários.
No Gráfico 4.27, é possível observar que há diferenças estatisticamente significativas entre CEFET/IF e duas Organizações Acadêmicas : Faculdades e Universidades ( ver também Tabela II.14, no Anexo II ).
Em relação às Categorias Administrativas, o Desconhecimento do conteúdo foi a opção escolhida por 19,7 \% dos estudantes, nas IES Públicas, e por 20,8 \%, nas IES Privadas.
Não são observadas diferenças estatisticamente significativas entre essas duas categorias.
A alternativa modal entre os alunos, quando agregados pelas Categorias Administrativas, para a dificuldade encontrada, foi causada pela Forma diferente de abordagem do conteúdo : 38,3 \%, entre os estudantes de IES Públicas, e 42,0 \%, entre os de IES Privadas ( ver Gráfico 4.28 e Tabela II.14, no Anexo II ). 4.6 CONTEÚDOS DAS QUESTÕES OBJETIVAS DA PROVA Ao se analisarem os conteúdos das questões objetivas da prova ( Questão 8 ), verifica -se que um percentual muito pequeno dos estudantes avaliados, apenas 6,9 \%, afirmou que não estudou ainda a maioria desses conteúdos ( Gráficos 4.29 a Gráfico 4.32, e as Tabelas II.15 e II.16, no Anexo II ).
A maior parte ( 47,0 \% ) afirmou ter estudado e aprendido muitos ou todos os conteúdos avaliados.
Na análise por Grande Região, a proporção de respondentes que escolheram a opção não estudou ainda a maioria desses conteúdos foi pequena.
Observa -se que, na região Sudeste ( 7,2 \% ), a proporção foi ligeiramente maior do que a da média nacional ( 6,9 \% ).
Nas regiões Nordeste, Sul e Centro-Oeste, a situação foi \textbf{inversa}, com proporções mais baixas, correspondendo a, respectivamente, 5,6 \%, 6,7 \% e 6,7 \%.
Na região Norte, a proporção de respondentes que escolheram essa resposta foi a mesma da nacional ( 6,9 \% ).
Não são observadas diferenças estatisticamente significativas entre as Grande Regiões.
Em todas as regiões, a maior parte dos presentes afirmou ter estudado e aprendido muitos ou todos os conteúdos, com proporções variando entre 37,8 \%, na região Norte, a 55,4 \%, na região Nordeste.
Considerando -se separadamente as opiniões de estudantes dos quatro quartos de desempenho, observa -se que, no quarto inferior, 11,7 \% ofereceram como resposta que não estudou ainda a maioria desses conteúdos, sendo 2,7 \% os do quarto superior com a mesma resposta.
As diferenças entre os alunos que optaram por esse motivo de dificuldade entre todos os quartos de desempenho são estatisticamente significativas.
Tendo -se em conta o quarto superior, 63,8 \% dos alunos afirmaram ter estudado e aprendido muitos ou todos os conteúdos.
No outro extremo, no primeiro quarto, 35,4 \% optaram pelas mesmas categorias.
As proporções são crescentes com o desempenho.
Na análise por Organização Acadêmica, a proporção de respondentes que escolheram a opção não estudou ainda a maioria desses conteúdos foi pequena.
Observa -se que as proporções nos CEFET/IF e Faculdades, 4,8 \% e 5,4 \%, respectivamente, foram menores do que a média nacional ( 6,9 \% ).
Nos Centros Universitários e Universidades, a situação foi \textbf{inversa} : a proporção foi mais alta ( 7,2 \% e 8,0 \%, respectivamente ).
No Gráfico 4.31, é possível observar que as diferenças de Faculdades e CEFET/IF em relação aos demais tipos de Organização Acadêmica ( Centros Universitários e Universidades ) são estatisticamente significativas, mas não entre si.
Em todos os tipos de Organização Acadêmica, a maioria dos presentes afirmou ter estudado e aprendido muitos ou todos os conteúdos, com proporções variando entre 42,0 \%, nas Universidades, a 60,5 \%, nos CEFET/IF ( ver também Tabela II.16, no Anexo II ).
Em relação às Categorias Administrativas, observa -se que, nas IES Públicas, 4,4 \% ofereceram como resposta que não estudou ainda a maioria desses conteúdos, sendo 7,9 \% nas IES Privadas com a mesma resposta.
A diferença entre os alunos que optaram por esse motivo de dificuldade nas duas categorias é estatisticamente significativa.
Considerando -se as IES Públicas, 58,3 \% dos alunos afirmaram ter estudado e aprendido muitos ou todos os conteúdos, sendo de 42,9 \% a proporção dos que escolheram essa resposta nas IES Privadas ( ver Gráfico 4.32 e Tabela II.16, no Anexo II ). 4.7 TEMPO GASTO PARA CONCLUIR A PROVA Ao responderem sobre o tempo de conclusão da prova ( Questão 9 ), quase dois terços dos estudantes ( 64,2 \% ) afirmaram ter gasto Entre duas e quatro horas ( Gráfico 4.33 a Gráfico 4.36 e, no Anexo II, as Tabelas II.17 e II.18 ).
Considerando -se as cinco Grandes Regiões brasileiras, os percentuais dos que utilizaram Entre duas e quatro horas para finalizar a prova nas regiões Sudeste ( 63,2 \% ) e Sul ( 64,1 \% ) foram ligeiramente inferiores ao percentual nacional.
Nas demais Grandes Regiões, os percentuais de alunos que \textbf{dispensaram} Entre duas e quatro horas para concluir a prova foram mais elevados que a média : 72,0 \% na região Norte ; 69,4 \%, na região Nordeste ; e 66,1 \%, na região Centro-Oeste, como mostrado no Gráfico 4.33.
Há diferenças estatisticamente significativas da região Norte e da região Nordeste em relação às regiões Sudeste e Sul, mas não entre si.
Uma vez analisadas as alternativas escolhidas pelos estudantes que se situam nos diferentes quartos de desempenho, observa -se uma tendência crescente : uma proporção maior de participantes no quarto superior declarou ter gasto Entre duas e quatro horas para concluir a prova quando comparada com a proporção de participantes nos quartos inferiores.
Há diferenças estatisticamente significativas entre todos os quartos de desempenho.
Na análise por Organização Acadêmica, os percentuais dos que utilizaram Entre duas e quatro horas para finalizar a prova nas Faculdades ( 65,3 \% ), Centros Universitários ( 65,7 \% ) e CEFET/IF ( 69,4 \% ) foram superiores ao percentual nacional ( 64,2 \% ).
Nas Universidades, tal percentual foi de 61,8 \%, e, como mostrado no Gráfico 4.35, há diferenças estatisticamente significativas entre Universidades e as proporções dos demais tipos de Organização Acadêmica, bem como entre Faculdades e CEFET/IF.
Em relação às Categorias Administrativas, observa -se que, nas IES Públicas, 66,3 \% declararam ter gasto Entre duas e quatro horas para concluir a prova, sendo 63,5 \% nas IES Privadas com a mesma resposta.
A diferença entre os alunos que optaram por esse motivo de dificuldade nas duas categorias é estatisticamente significativa ( ver Gráfico 4.36 e Tabela II.18, no Anexo II ).

% matched lemmas: arredondamento, crescer, disjunto, dispensar, estimador, estudantes, indagar, interseção, inverso, numérico, ordenar, parâmetro, percentagem, questionário, recorte, reprodução, subconjunto, sumário, totalizar, vidar
\end{textsample}
